
This paper evaluates ExtrospectiveNLI, a post-hoc hallucination detection approach that applies frozen NLI contradiction scoring to pre-existing (context, response) pairs, requiring no LLM generation at inference time. Using \texttt{cross-encoder/nli-deberta-v3-large} (304M parameters) on HaluEval across 60,000 balanced pairs (20,000 per task), the approach yields AUROC = 0.709 on dialogue and AUROC = 0.644 on QA. On summarization, AUROC = 0.530, near chance. Against SelfCheckGPT-NLI evaluated under the same benchmark, these figures represent gains of +0.229 and +0.114 respectively on the two commission-type tasks. Mechanistic analysis via Wilcoxon rank-sum tests (p $\leq$ 2.$07\times 10^{-13}$ on all three tasks; p $\approx$ 0 for dialogue) and Cohen's d (0.714 and 0.779 for dialogue and QA) confirms that DeBERTa's contradiction scores are stochastically ordered by hallucination label across all tasks. The near-chance result on summarization is attributed to the commission/omission distinction in hallucination types: NLI contradiction scoring detects commission-type hallucinations (fabricated facts, factual substitutions) but is architecturally unsuited to detecting omission-type hallucinations (abstractive gaps, missing information) that predominate in HaluEval-Summarization. The theoretical maximum AUROC for contradiction-based detection on summarization is approximately 0.52 given p\_contradictory $\approx$ 0.04 per example; the observed 0.530 is near this ceiling. Ablation experiments evaluating individual design choices (net-contradiction framing, sentence-level aggregation, dialogue windowing) were not executed, constituting a recognized limitation.
